\section{Problem formulation}
\label{sec:problem}

\paragraph{Setting.} An agent answers a question about the current (or a historical) value of an attribute by
calling a memory-search tool and then submitting a final answer. The store contains a group of records about that
attribute, written at different times, some of which assert values that disagree. The agent's terminal action is
three-way: submit a value, submit \cond{conflict} when the evidence does not determine one, or abstain.

\paragraph{What a resolution mechanism must decide.} Given a group of records $R=\{r_1,\dots,r_k\}$ about one
attribute, a mechanism must produce, for each record, a status (established, superseded, historical, unresolved)
and --- this is the part we vary --- must decide what to \emph{do} with that status at the terminal action. Three
designs span the space we study:

\begin{description}[leftmargin=1.2em]
\item[Advisory.] The status is shown to the agent and nothing more. The agent may ignore it. Conditions
\cond{G} (metadata labels), \cond{T} (temporal labels), \cond{EB} (evidence-basis labels) and \cond{LA} (LLM
arbiter) are advisory.
\item[Ungated enforcement.] Every resolution the mechanism produces is enforced at the terminal action: a
submission that disagrees is blocked or corrected. Conditions \cond{R}, \cond{RT} and \cond{EBG-all} enforce
without discrimination.
\item[Gated enforcement.] Only resolutions whose evidence basis is strong enough are enforced; the rest are
advisory. Conditions \cond{EBG} (and its variants) and \cond{LG} are gated; they differ in \emph{who} produces the
resolutions.
\end{description}

\paragraph{Why the distinction matters.} A reliability label is itself an inference, and it can be wrong in two
ways that have opposite remedies. It can be \emph{under-determined} --- the metadata cannot order two records, so
the mechanism reports a conflict that the text would resolve --- and it can be \emph{over-determined} --- the
mechanism asserts an ordering the evidence does not support. Advisory labelling leaves both errors to the agent,
which for small models means they are often followed anyway. Ungated enforcement removes the agent's discretion,
which converts an over-determined label directly into a wrong committed answer. Gating is the proposal that a
mechanism should enforce only where it can say \emph{why} it is sure.

\paragraph{Outcomes.} The frozen scorer classifies each run as correct or into one of several failure categories.
Two matter throughout. An answer is \textbf{harmful} if it commits to a wrong value: \cond{wrong\_value},
\cond{stale\_value}, \cond{unsupported\_claim} or \cond{arbitrary\_pick}. A group that genuinely cannot be ordered
is \textbf{falsely resolved} if the agent asserts a value for it instead of reporting a conflict. Task success and
harm are reported separately throughout, because, as \Cref{sec:ablations} shows, they do not move together.

\paragraph{Unit of analysis.} The unit is the task. Success per task is the mean over a model's decodings; contrasts
are paired within model and decoding; the clustering unit is the task \emph{family}, a (category, attribute) pair.
This matters because tasks within a family share a template and are not independent.
